\section{Threats to Validity}
\label{sec:validity}

\paragraph{Construct validity.}
The main construct is reproduction of a specified vulnerability in an
affected original target. Product names, versions, component boundaries,
entry points, and effects can each be misidentified. Reference profiles
assembled from public information can propagate incorrect assumptions.
An original component tested through a bypassed frontend can also be
mistaken for reproduction of the complete documented attack. The proposed
evidence dimensions make these assumptions explicit, but the per-case
audit and replay remain unfinished. \todo{VAL-01}{Report the actual number
of independently checked targets, replayed outcomes, available controls,
and unresolved disagreements; align the claims with that coverage.}

Behavior labels introduce another construct risk. Describing a saved
script as a simulation is often possible from source; identifying the
reason the agent chose it usually needs a trajectory. A visible error
does not directly reveal the model's reasoning, and labels such as
``ignored feedback'' require observable opportunities to respond.
LLM-assisted retrieval can aid inspection but does not make a second
agent's judgment independent ground truth. Human coding procedures and
ambiguous cases must therefore be reported as actually performed.

\paragraph{Internal validity.}
Observed performance combines the model route, harness, tools, task
wording, host, budget, and external services. The single ARM64 host can
favor native execution of some targets, while other cases require
cross-architecture emulation. Independent containers were intended to
limit carryover, but shared network state, caches, package sources, and
provider availability may remain. Model snapshots and serving behavior
may change during a collection period. The unfinished configuration audit
limits attribution of observed differences.

Missing traces are particularly consequential for process claims.
Excluding them could bias response distributions if missingness is
associated with a configuration or difficult case. Including a historical
label without its primary evidence would conceal that uncertainty. We
will report coverage by analysis and distinguish recording loss from
known execution failures where possible. Unknown reasons will remain
unknown. No causal claim about the benefit of an observed repair or the
effect of an obstacle follows from the natural choice of actions.

\paragraph{External validity.}
The 30 cases were selected for public documentation, apparent firmware
accessibility, and three vulnerability classes. This favors cases that
can be investigated from public sources and excludes much of IoT's
hardware, protocol, and vulnerability diversity. Equal class sizes do
not represent the natural prevalence of weaknesses. Multiple CVEs may
share a vendor, component, image, or dependency, further reducing effective
diversity. Findings should be framed around the studied situations,
not the overall rate at which arbitrary IoT CVEs can be reproduced.

Public CVE descriptions and PoCs may also have appeared in model training
or accessible retrieval sources. The study evaluates execution with
public information available; it does not isolate reasoning about unseen
vulnerability mechanisms. Results from one harness and budget do not
establish performance under another agent architecture or later model
snapshot. \todo{EXP-03}{Optional, if a principal finding is claimed to
generalize across harnesses: replicate a preselected, diverse subset under
a second harness with the same model route and comparable resources.}

\paragraph{Conclusion validity.}
The 450 positions are repeated observations on 30 cases, not 450
independent vulnerabilities. Three trials offer limited resolution for
consistency, and small class or ISA subgroups cannot support broad
comparative claims. Paired case comparisons and uncertainty should reflect
clustering and the specific estimand. Exploratory subgroup analyses and
multiple patterns should be presented as exploratory, including null or
contradictory observations.

The inspected examples were chosen because they illuminate different
evidence boundaries. They do not establish how often those patterns
occur. Cross-case coding, transparent case-selection criteria, and
negative examples are needed to support stronger findings.
\todo{VAL-02}{Report final coverage, sampling, counterexamples, missing-data
sensitivity, and any exploratory comparisons omitted for insufficient
evidence. Remove or narrow interpretations that do not survive these
checks.}
